\section{Supporting Study: Graph Learning and Bilinear Gating}
\label{sec:certgraph_arch}

\setcounter{theorem}{0}
\renewcommand{\thetheorem}{A.\arabic{theorem}}
\providecommand{\theHtheorem}{}
\renewcommand{\theHtheorem}{appendixA.\arabic{theorem}}
\renewcommand{\theproposition}{A.\arabic{proposition}}

This supporting study examines whether learned graph representations help identify the stated resolved-permission predicate (Figure~\ref{fig:certgraph_arch}). Certificate evidence is collected independently of its predictions. The conditional analysis below motivates matched experiments; it does not establish a failure of every attention model. The evaluated bilinear scorer provides no classification advantage over additive attention, so the certificate-lineage contribution does not depend on it.

\subsection{Heterogeneous Graph Attention Network (HetGAT)}
Following graph-learning formulations~\cite{scarselli2008graph,wu2021survey}, CertGraph employs a multi-relational graph attention network. Let $\mathbf{x}_v \in \mathbb{R}^{d_{\phi(v)}}$ denote the initial feature vector of entity $v \in \mathcal{V}$. Vertices are first projected into a shared hidden space of dimension $d$:
\begin{equation}
\fitmath{\mathbf{h}_v^{(0)} = \text{ReLU} \left( \mathbf{W}_{\phi(v)}^{(0)} \mathbf{x}_v + \mathbf{b}_{\phi(v)}^{(0)} \right)}
\end{equation}
where $\mathbf{W}_{\phi(v)}^{(0)} \in \mathbb{R}^{d \times d_{\phi(v)}}$ is a type-specific linear projection matrix.

At each layer $l \in \{1, \dots, L\}$, message passing aggregates relational representations across incoming edges $\mathcal{N}_r(v) = \{u \in \mathcal{V} \mid (u, v) \in \mathcal{E}, \psi(u, v) = r\}$:
\begin{equation}
\fitmath{\mathbf{h}_v^{(l)} = \sigma \left( \sum_{r \in \mathcal{R}_{\text{in}}(\phi(v))} \sum_{u \in \mathcal{N}_r(v)} \alpha_{vu}^{(l)} \mathbf{W}_r^{(l)} \mathbf{h}_u^{(l-1)} \right)}
\end{equation}
where $\mathbf{W}_r^{(l)} \in \mathbb{R}^{d \times d}$ is the relation-specific weight transformation for relation $r \in \mathcal{T}_E$, and $\sigma(\cdot)$ denotes the ELU activation function. The relational attention coefficient $\alpha_{vu}^{(l)}$ assigns dynamic importance to edge $(u, v)$ under relation $r$:
\begin{equation}
\fitmath{\begin{gathered}
\alpha_{vu}^{(l)}=\frac{\exp(z_{vu}^{(l)})}{\sum_{k\in\mathcal N_r(v)}\exp(z_{vk}^{(l)})},\\
z_{vu}^{(l)}=\operatorname{LeakyReLU}\!\left(\mathbf a_r^\top
\begin{bmatrix}\mathbf W_r^{(l)}\mathbf h_v^{(l-1)}\\ \mathbf W_r^{(l)}\mathbf h_u^{(l-1)}\end{bmatrix}\right).
\end{gathered}}
\end{equation}
where $\mathbf{a}_r \in \mathbb{R}^{2d}$ is the parameterized attention vector for relation $r$, and $\|$ denotes vector concatenation.

\begin{figure*}[ht]
\centering
\includegraphics[width=\textwidth]{figures/certgraph_learning_architecture.pdf}
\caption{CertGraph learning architecture and the literal bilinear mechanism. (A) Typed graph features enter (B) two relational attention layers with four heads per layer and gated self-skips, followed by (C) template readout and binary classification in the matched resolved-permission ESC1 study. (D) Source and target states determine a scalar bilinear gate that scales the transformed neighbor message before attention-weighted aggregation. Exact authorization is checked separately; the learned gate is not a Boolean reachability test. The input topology is illustrative.}
\label{fig:certgraph_arch}
\end{figure*}

\subsection{Theoretical Analysis: Representation Collapse and Gradient Bounds}
One limitation in pure neighborhood aggregation on directed multigraphs is the erasure of intrinsic attributes for leaf and intermediate nodes.

\begin{theorem}[Intrinsic Attribute Preservation and Gradient Lower Bounds]
\label{thm:theorem1}
Let $\mathcal{G} = (\mathcal{V}, \mathcal{E}, \phi, \psi)$ be a directed heterogeneous multigraph. Consider an $L$-layer pure relational aggregation model with no target-feature dependence in its aggregation coefficients, no biases, and $\sigma(\mathbf{0})=\mathbf{0}$. For the residual bound, differentiate with respect to the common hidden input $\mathbf h_v^{(0)}$ and assume square compatible skip maps and differentiability. No lower bound for the raw type-projected features is asserted. These assumptions restrict the theorem; attention coefficients that depend on the target features can transmit those features even without a self-loop.
\begin{enumerate}
    \item \textbf{Intrinsic Feature Erasure:} If $\mathcal{G}$ contains no explicit self-loops on template nodes ($v \notin \mathcal{N}_r(v)$ for all $r$), the first-layer representation $\mathbf{h}_v^{(1)}$ is conditionally independent of initial features $\mathbf{x}_v$:
    \begin{equation}
\fitmath{\frac{\partial \mathbf{h}_v^{(1)}}{\partial \mathbf{x}_v} = \mathbf{0} \in \mathbb{R}^{d \times d_{\phi(v)}}}
\end{equation}
    Furthermore, for any unprivileged source account $u$ with in-degree $d_{\text{in}}(u) = 0$, $\mathbf{h}_u^{(l)} = \mathbf{0}$ for all $l \ge 1$.
    \item \textbf{Guaranteed Gradient Lower Bound via Residual Skips:} When parameterized residual skip connections are introduced ($\tilde{\mathbf{h}}_v^{(l)} = \mathbf{h}_v^{(l)} + \mathbf{W}_{\text{skip}}^{(l)} \tilde{\mathbf{h}}_v^{(l-1)}$) with $\sigma_{\min}(\mathbf{W}_{\text{skip}}^{(l)}) > 0$, the Jacobian of the representation satisfies:
    \begin{equation}
\fitmath{\sigma_{\min} \left( \frac{\partial \tilde{\mathbf{h}}_v^{(L)}}{\partial \mathbf{h}_v^{(0)}} \right) \ge \prod_{k=1}^L \sigma_{\min}(\mathbf{W}_{\text{skip}}^{(k)}) - \|\mathcal{J}_{\text{graph}}(v)\|}
\end{equation}
    where $\mathcal{J}_{\text{graph}}(v)$ is the full derivative of all non-skip contributions, including their dependence through earlier residual states. Its norm must be measured or bounded for the particular model; a single path's weight product does not bound an arbitrary graph's aggregate derivative. When $\prod_{k=1}^L \sigma_{\min}(\mathbf{W}_{\text{skip}}^{(k)}) > \|\mathcal{J}_{\text{graph}}(v)\|$, initial feature sensitivity is strictly preserved:
    \begin{equation}
\fitmath{\left\| \frac{\partial \tilde{\mathbf{h}}_v^{(L)}}{\partial \mathbf{h}_v^{(0)}} \right\| > 0}
\end{equation}
\end{enumerate}
\end{theorem}

\begin{proof}
We establish the two claims sequentially:

\paragraph{Part 1: Intrinsic Feature Erasure under Pure Neighborhood Aggregation.}
Under pure relational convolution at layer $l=1$, the hidden state $\mathbf{h}_v^{(1)}$ is computed as $\sigma \left( \sum_{r \in \mathcal{R}_{\text{in}}(\phi(v))} \sum_{u \in \mathcal{N}_r(v)} \alpha_{vu}^{(1)} \mathbf{W}_r^{(1)} \mathbf{h}_u^{(0)} \right)$. In an Active Directory multigraph, template nodes possess no explicit self-loops ($v \notin \mathcal{N}_r(v)$ for all $r$). Differentiating $\mathbf{h}_v^{(1)}$ with respect to the initial template feature vector $\mathbf{x}_v$ through its type-specific initial projection:
\begin{equation}
\fitmath{\begin{gathered}\frac{\partial \mathbf{h}_v^{(1)}}{\partial \mathbf{x}_v} \\ = \sigma'(\cdot) \sum_{r \in \mathcal{R}_{\text{in}}(\phi(v))} \sum_{u \in \mathcal{N}_r(v)} \alpha_{vu}^{(1)} \mathbf{W}_r^{(1)} \frac{\partial \mathbf{h}_u^{(0)}}{\partial \mathbf{x}_v}\end{gathered}}
\end{equation}
Because distinct nodes possess independent initial feature vectors ($\partial \mathbf{h}_u^{(0)} / \partial \mathbf{x}_v = \mathbf{0}$ for all $u \neq v$), every term in the summation vanishes identically, yielding $\partial \mathbf{h}_v^{(1)} / \partial \mathbf{x}_v = \mathbf{0}$. Thus, the layer-1 embedding is conditionally independent of $\mathbf{x}_v$, discarding all intrinsic template configuration flags (e.g., \texttt{CT\_FLAG\_ENROLLEE\_SUPPLIES\_SUBJECT}). Moreover, for any source node $u$ with in-degree $d_{\text{in}}(u) = 0$, $\mathcal{N}_r(u) = \emptyset$ for all $r$, so $\mathbf{h}_u^{(l)} = \mathbf{0}$ for all $l \ge 1$ by empty summation and $\sigma(\mathbf{0}) = \mathbf{0}$.

\paragraph{Part 2: Gradient Lower Bounds via Parameterized Residual Skips.}
With parameterized residual skip connections $\tilde{\mathbf{h}}_v^{(l)} = \mathbf{h}_v^{(l)} + \mathbf{W}_{\text{skip}}^{(l)} \tilde{\mathbf{h}}_v^{(l-1)}$ and $\tilde{\mathbf{h}}_v^{(0)} = \mathbf{h}_v^{(0)}$, unrolling the recurrence across $L$ layers yields:
\begin{equation}
\fitmath{\begin{gathered}\tilde{\mathbf{h}}_v^{(L)} = \left( \prod_{k=1}^L \mathbf{W}_{\text{skip}}^{(k)} \right) \mathbf{h}_v^{(0)} \\ + \sum_{l=1}^L \left( \prod_{k=l+1}^L \mathbf{W}_{\text{skip}}^{(k)} \right) \mathbf{h}_v^{(l)}\end{gathered}}
\end{equation}
Differentiating $\tilde{\mathbf{h}}_v^{(L)}$ with respect to $\mathbf{h}_v^{(0)}$:
\begin{equation}
\fitmath{\begin{gathered}\frac{\partial \tilde{\mathbf{h}}_v^{(L)}}{\partial \mathbf{h}_v^{(0)}} = \prod_{k=1}^L \mathbf{W}_{\text{skip}}^{(k)} + \mathcal{J}_{\text{graph}}(v),\\ \text{where}\quad \mathcal{J}_{\text{graph}}(v) = \sum_{l=1}^L \left( \prod_{k=l+1}^L \mathbf{W}_{\text{skip}}^{(k)} \right) \frac{\partial \mathbf{h}_v^{(l)}}{\partial \mathbf{h}_v^{(0)}}\end{gathered}}
\end{equation}
By Weyl's singular value perturbation inequality and the sub-multiplicative property of singular values:
\begin{equation}
\fitmath{\begin{gathered}\sigma_{\min} \left( \frac{\partial \tilde{\mathbf{h}}_v^{(L)}}{\partial \mathbf{h}_v^{(0)}} \right) \\ \ge \sigma_{\min} \left( \prod_{k=1}^L \mathbf{W}_{\text{skip}}^{(k)} \right) - \|\mathcal{J}_{\text{graph}}(v)\| \\ \ge \prod_{k=1}^L \sigma_{\min}(\mathbf{W}_{\text{skip}}^{(k)}) - \|\mathcal{J}_{\text{graph}}(v)\|\end{gathered}}
\end{equation}
Under the stated coefficient-independence assumptions, absence of every directed return walk to $v$ within the evaluated depth implies $\mathcal{J}_{\text{graph}}(v)=\mathbf{0}$. Then the minimum singular value of the skip product is bounded below by the product of the individual minimum singular values; equality is not generally implied. In cyclic or feature-dependent attention models, the non-skip derivative can reinforce or cancel the skip path. The positive-sensitivity conclusion therefore requires the displayed strict inequality and is not an unconditional guarantee for all GNNs.
\end{proof}

\subsection{Shortcut Behavior and Bilinear Gating}
While Theorem~\ref{thm:theorem1} gives a conditional sensitivity bound for additive residual skips ($\tilde{\mathbf{h}}_v^{(l)} = \mathbf{h}_v^{(l)} + \mathbf{W}_{\text{skip}}^{(l)} \tilde{\mathbf{h}}_v^{(l-1)}$), empirical testing reveals a secondary failure mode: \textit{inductive shortcut collapse}.

\paragraph{The Shortcut Memorization Failure}
Because certificate template attributes (e.g., \texttt{CT\_FLAG\_ENROLLEE\_SUPPLIES\_SUBJECT} $= 1$ and Client Authentication EKU $= 1$) are strongly correlated with positive escalation labels in synthetic training distributions, additive skip connections allow the network to bypass the multi-hop graph aggregation entirely. The MLP classification head learns an inductive shortcut:
\begin{equation}
\fitmath{\hat{y}_t = \text{MLP}(\tilde{\mathbf{h}}_t^{(L)}) \approx \text{MLP} \left( \prod_{k=1}^L \mathbf{W}_{\text{skip}}^{(k)} \mathbf{h}_t^{(0)} \right)}
\end{equation}
When the trained model is evaluated against \textbf{adversarial hard negatives}—templates that possess vulnerable attribute flags but reside in disconnected graph components or require mandatory manager approval—standard additive-skip GNNs memorize the flag pattern and classify them as vulnerable, showing poor agreement on the historical canonical generator. The matched audit in Section~\ref{sec:evaluation} corrects its reachable-path labels and separates augmentation from implementation differences.

\paragraph{Multiplicative Bilinear Reachability Gating (\texttt{CertGraph-Gated}):}
To reduce shortcut memorization, we condition messages on a learned bilinear compatibility gate. The division between learned scoring and symbolic checking follows a broader neuro-symbolic design~\cite{garcez2023neurosymbolic}. Exact reachability is checked separately by the existing symbolic verification step; the learned gate alone is not a Boolean authorization predicate. We formulate the \textbf{Multiplicative Bilinear Compatibility Gate}:
\begin{equation}
\fitmath{\mathbf{g}_{uv}^{(l)} = \sigma \left( \mathbf{h}_u^{(l-1)\top} \mathbf{W}_g^{(l)} \mathbf{h}_v^{(l-1)} + b_g^{(l)} \right)}
\end{equation}
where $\mathbf{W}_g^{(l)} \in \mathbb{R}^{d \times d}$ is a bilinear interaction matrix, $b_g^{(l)} \in \mathbb{R}$ is a scalar bias, and $\sigma(\cdot)$ is the sigmoid function mapping to $(0, 1)$.

The gated message vector $\mathbf{m}_{vu}^{(l)}$ and state update are defined as:
\begin{equation}
\fitmath{\begin{gathered}\mathbf{m}_{vu}^{(l)} = \mathbf{g}_{uv}^{(l)} \odot \left( \mathbf{W}_r^{(l)} \mathbf{h}_u^{(l-1)} \right)\end{gathered}}
\end{equation}
\begin{equation}
\fitmath{\begin{gathered}\mathbf{h}_v^{(l)} = \sigma \left( \sum_{r \in \mathcal{R}_{\text{in}}(\phi(v))} \sum_{u \in \mathcal{N}_r(v)} \alpha_{vu}^{(l)} \mathbf{m}_{vu}^{(l)} \right)\\ \quad+\mathbf{g}_{vv}^{(l)} \odot \left( \mathbf{W}_{\text{skip}}^{(l)} \mathbf{h}_v^{(l-1)} \right)\end{gathered}}
\end{equation}
where $\odot$ denotes the element-wise Hadamard product.

\begin{proposition}[Conditional Attenuation by Multiplicative Gating]
For a bounded transformed neighbor message, if its scalar gate tends to zero, multiplicative attenuation gives:
\begin{equation}
\fitmath{\lim_{\mathbf{g}_{uv} \to 0} \mathbf{m}_{vu}^{(l)} = \mathbf{0}}
\end{equation}
Disconnection alone does not imply this gate limit: with zero embeddings and zero bias, the sigmoid gate equals $1/2$. Other residual paths or classifier biases can also produce a positive output. The attenuation statement therefore does not prove universal elimination of shortcuts. The historical numerical implementation uses a feature-channel gate rather than this literal scorer. We now implement a separate two-layer, four-head realization of the displayed bilinear edge and self-skip equations and evaluate it with identical-initialization additive controls and matched augmentation (Table~\ref{tab:literal_bilinear}). All relevant parameter families are trainable through autograd. The two implementations have identical held-out classifications in that new study, and both miss the positive long-path cases. This validates an auditable realization of the formulation, rather than a universal shortcut-elimination or authorization guarantee. A second-tier verifier must process unresolved and out-of-horizon cases, including neural negatives, if it is to recover those observed misses.
\end{proposition}

\paragraph{Synthetic Identity Graphs} The historical static benchmark uses 700 class-balanced synthetic identity graphs. Additional graph sizes are used in the separate scalability experiments. Neither the class proportions nor the graph generator constitute a measured enterprise distribution. The new matched gating audit uses the existing two-step implementations with fixed source snapshots and separately recorded training, validation, and test inputs.

\subsection{In-Distribution Ablations, Calibration, and Perturbation Asymmetry}
We first distinguish the historical class-balanced generator from the newer attribution and live-policy experiments. Their labels, model implementations, and observation units differ.

\paragraph{Skip-Connection Ablation and Conditional Theory:}
The saved five-fold artifact reports Macro-F1 $0.9986 \pm 0.0029$ with skips and $0.4768 \pm 0.0239$ without skips. Recomputing the paired statistic from the five saved fold values gives $t=46.245$, $p=1.308\times10^{-6}$, and four degrees of freedom. The dispersions are population standard deviations across these folds. Shared training data and generator assumptions limit interpretation of this descriptive comparison. An empirical ablation does not validate the theorem outside its coefficient-independence and derivative assumptions.

\paragraph{Class Proportions and Probability Interpretation}
The static experiments reported here are generator-based classification studies. Their balanced labels do not establish a 96\% enterprise safe base rate, and their softmax scores are not calibrated probabilities of real-world compromise. We retain probabilities and selected checkpoints in the matched audit so that classification errors can be checked independently.

\paragraph{Perturbation Asymmetry and Simple Comparators}
The saved robustness artifact gives Macro-F1 $0.9479$ after 30\% edge perturbation, compared with $0.8965$ and $0.7709$ after 5\% and 10\% feature perturbation. This sensitivity pattern motivates testing structural counterfactuals; it does not prove that every prediction ignores topology. A separately verified, previously saved three-seed ESC13 abstraction also favors simple graph-count models over the existing GNN on its IID cases. Its symbolic reference defines the labels, so its perfect agreement is a consistency result rather than independent exploit validation.

\subsection{Static GNN Generalization and Adversarial Hard Negatives}
A source audit identifies two problems with interpreting the historical $41/63$ safe-label agreement as a general hard-negative detection result. Ten ``deep horizon'' cases retain a membership path to an allowed enrollment group and therefore are positive under a transitive effective-permission interpretation. The stored 22 disagreements are eight administrator-only write cases, four deny-write cases, seven benign OID links, and three deep OID links; the earlier OID/CA explanation did not match that artifact. The canonical generator also repeats inputs across training and validation. These historical counts are retained in the artifact but do not establish independent-case specificity or a binomial confidence interval.

We therefore freeze and execute a matched $2\times2$ diagnostic: additive versus feature-channel-gated implementations, each with counterfactual augmentation probability zero or 0.3, across three new seeds. Model equations, propagation depth, and original artifacts remain unchanged. Each seed evaluates 250 opposite-label families, with membership depths 0, 1, 2, 4, and 8, randomized node order, and matched configuration features. Positive members retain reachable enrollment; negative members either restrict enrollment to the administrator or add a reachable explicit deny. Queue traversal and Boolean closure independently agree on all test labels under this stated ESC1 permission abstraction. CA trust and resolved permissions are assumed; Windows exploitability is not the label definition.

% Generated from independently verified saved checkpoint predictions.
\begin{table}[tbp]
\centering
\caption{Matched two-step diagnostic: 1,500 paired cases per arm (750 positive, 750 negative). Correct classification is a percentage; false positives and negatives are counts. Repeated training archetypes limit this synthetic diagnostic.}
\label{tab:hard_negatives}
\evidencetable
\begin{tabular}{llrrr}
\toprule
Implementation & Aug. & Correct & \makecell{False\\positive} & \makecell{False\\negative} \\
\midrule
Additive & 0 & 50.00\% & 750 & 0 \\
Additive & 0.3 & 71.67\% & 425 & 0 \\
\makecell[l]{Feature-channel\\gate} & 0 & 70.00\% & 0 & 450 \\
\makecell[l]{Feature-channel\\gate} & 0.3 & 70.00\% & 0 & 450 \\
\bottomrule
\end{tabular}
\end{table}


The gate gives zero false positives but misses all 450 positive cases at depths 2, 4, and 8 across the three seeds. Adding the same augmentation to the additive implementation improves its pooled accuracy from 50\% to 71.67\%, compared with the gate's 70\%. Thus the audit does not establish superiority of gating over an augmentation-matched comparator. First-step membership transforms are fixed in these existing NumPy implementations, and their feature-channel gate is $\sigma(W_{\mathrm{gate}}\mu_{\mathrm{reach}}+b)$, rather than a literal implementation of the general bilinear edge scorer in Section~\ref{sec:certgraph_arch}. Initialization also differs between implementations. These distinctions prevent attributing the entire contrast to a bilinear operation.

Only 22 unique training inputs occur per seed, with 21 also present in validation. Consequently, validation selection in this diagnostic is not an independent-distribution generalization test. All 6,000 saved test predictions are reproduced from saved checkpoints, and their labels and aggregates are checked independently. A symbolic second tier must also examine unresolved or out-of-horizon cases: verifying only neural-positive candidates cannot recover these observed neural false negatives.

\subsubsection{Literal bilinear evaluation}
We address the implementation mismatch with a separate literal realization of the bilinear edge and self-skip equations: two relational-attention layers, four heads, 32 hidden channels, type-specific input projection, and fully trainable encoder, attention, skip, and gate parameters. The binary readout evaluates the stated ESC1 resolved-permission predicate, not observed Windows compromise. Additive and gated arms start with identical parameter values and receive matched augmentation. Neither symbolic closure nor outcome labels enter the forward pass. Gradient, node-relabeling, and padding checks validate implementation properties; the historical implementations remain unchanged.

The new generator constructs varied reachable and unreachable membership branches and paired allow/deny interventions. For each of three new seeds, 300 training and 100 validation families are separated from 150 families in each of IID, long-path, and larger-background test cohorts. Typed five-round Weisfeiler--Lehman input hashes and family IDs are disjoint across splits within a seed. This conservatively excludes isomorphic whole inputs under the encoded roles and edge types, without claiming disjoint local motifs. Long-path positive enrollment chains have membership depth 2, 4, or 8. The larger-background cohort adds isolated nodes to different generated families; it is not a matched causal estimate of graph-size effects or enterprise-scale generalization.

% Generated from verified literal bilinear checkpoint predictions.
\begin{table}[ht]
\centering
\caption{Literal two-layer attention comparison: mean accuracy (\%) and sample standard deviation across three seeds, with 150 paired families per test cohort per seed. Every arm starts with identical parameter values. Augmentation is matched; typed whole-graph hashes and family IDs are disjoint across splits within each seed.}
\label{tab:literal_bilinear}
\small
\begin{tabular}{llrrr}
\toprule
Model & Augmentation & IID & Long paths & Larger background \\
\midrule
Additive attention & 0 & $86.00 \pm 2.19$ & $50.00 \pm 0.00$ & $87.00 \pm 1.00$ \\
Additive attention & 0.3 & $86.00 \pm 2.19$ & $50.00 \pm 0.00$ & $87.00 \pm 1.00$ \\
Bilinear attention & 0 & $86.00 \pm 2.19$ & $50.00 \pm 0.00$ & $87.00 \pm 1.00$ \\
Bilinear attention & 0.3 & $86.00 \pm 2.19$ & $50.00 \pm 0.00$ & $87.00 \pm 1.00$ \\
\bottomrule
\end{tabular}
\end{table}


Independent NetworkX reachability checks all 5,100 saved case labels, and checkpoint inference reproduces all 10,800 held-out predictions. The literal bilinear and additive models have identical classifications in every matched arm: IID accuracy averages $86.00\%\pm2.19\%$, larger-background accuracy $87.00\%\pm1.00\%$, and long-path accuracy 50\%. In each arm, all 450 positive long-path cases are missed. Some IID negatives also require deny-membership reasoning beyond the local horizon; 126 pooled false positives remain. Probability scores differ, but their reported calibration errors concern these synthetic cohorts only. This closes the literal-implementation and repeated-whole-input gaps while failing to establish classification superiority for bilinear gating. Exact verification of unresolved cases remains necessary under the stated permission contract.

